% =============================================================================
% preamble.tex
% Shared package configuration for the PhD proposal.
% This file is \input{} into proposal.tex — do NOT add \documentclass here.
% =============================================================================


% -----------------------------------------------------------------------------
% 1. ENCODING & FONTS
% Must be loaded first so all other packages interpret characters correctly.
% -----------------------------------------------------------------------------
\usepackage[T1]{fontenc}      % Use 8-bit font encoding (handles accents, etc.)
\usepackage[utf8]{inputenc}   % Accept UTF-8 characters in source files
\usepackage{lmodern}          % Latin Modern fonts — scalable, good for PDF
\usepackage{microtype}        % Subtle improvements to text justification & spacing


% -----------------------------------------------------------------------------
% 2. PAGE LAYOUT
% Controls paper size, margins, and line spacing.
% -----------------------------------------------------------------------------
\usepackage[a4paper, margin=2.5cm]{geometry}   % A4 paper with 2.5 cm margins
\usepackage{setspace}
\onehalfspacing                                 % 1.5 line spacing (common for theses)
                                                % Use \doublespacing if required by your university


% -----------------------------------------------------------------------------
% 3. MATHEMATICS
% Essential for a technical/scientific PhD proposal.
% -----------------------------------------------------------------------------
\usepackage{amsmath}      % Core math environments: align, equation, gather, etc.
\usepackage{amssymb}      % Extra math symbols: \mathbb{R}, \nabla, etc.
\usepackage{amsthm}       % Theorem, lemma, proof environments
\usepackage{mathtools}    % Extensions to amsmath (e.g. \coloneqq, dcases)
\usepackage{subfiles}     % Allows each chapter to be compiled separately (for faster editing)

% -----------------------------------------------------------------------------
% 4. FIGURES & GRAPHICS
% Tells LaTeX where to find images and how to handle them.
% -----------------------------------------------------------------------------
\usepackage{graphicx}     % \includegraphics command
\graphicspath{            % Search these folders for figures automatically
  {continuum_model/figures/}
  {molecular_dynamics/figures/}
  {sensitivity_analysis/figures/}
}
\usepackage{float}        % [H] placement option — forces figure exactly here
\usepackage{caption}      % Customise caption style and spacing
\usepackage{subcaption}   % \begin{subfigure} for side-by-side figures


% -----------------------------------------------------------------------------
% 5. TABLES
% Packages for clean, publication-quality tables.
% -----------------------------------------------------------------------------
\usepackage{booktabs}     % \toprule, \midrule, \bottomrule — professional rules
\usepackage{array}        % Extended column types in tabular (e.g. >{...}c)
\usepackage{longtable}    % Tables that span multiple pages


% -----------------------------------------------------------------------------
% 6. UNITS & NUMBERS
% Consistent formatting of physical quantities and units.
% Usage: \SI{300}{\kelvin}  \SI{1.5e-9}{\metre}  \num{1234567}
% -----------------------------------------------------------------------------
\usepackage{siunitx}
\sisetup{
  separate-uncertainty = true,    % e.g. 1.23 ± 0.04
  range-phrase = { to },          % e.g. 1 to 10 nm
}


% -----------------------------------------------------------------------------
% 7. CODE LISTINGS
% For displaying simulation scripts, input files, or pseudocode.
% -----------------------------------------------------------------------------
\usepackage{xcolor}       % Colour support (needed by listings for highlighting)
\usepackage{listings}     % \begin{lstlisting} environment
\lstset{
  basicstyle  = \ttfamily\small,
  keywordstyle= \color{blue},
  commentstyle= \color{gray},
  numbers     = left,
  numberstyle = \tiny\color{gray},
  breaklines  = true,
  frame       = single,
}


% -----------------------------------------------------------------------------
% 8. MISCELLANEOUS UTILITIES
% Small but useful packages for writing and drafting.
% -----------------------------------------------------------------------------
\usepackage{enumitem}     % Customise list environments (\begin{itemize}, etc.)
\usepackage{todonotes}    % \todo{fix this} — inline notes while drafting
                          % Use \listoftodos to see all todos at once


% -----------------------------------------------------------------------------
% 9. BIBLIOGRAPHY  (BibLaTeX + Biber backend)
% Single .bib file at the proposal/ root. Biber handles Unicode correctly.
% Run order: pdflatex → biber → pdflatex → pdflatex  (latexmk handles this)
% -----------------------------------------------------------------------------
\usepackage[
  backend  = biber,
  style    = authoryear,  % Change to 'numeric' if your field uses [1] citations
  sorting  = nyt,         % Sort by name, year, title
  maxcitenames = 2,       % "Smith & Jones" — more authors become "Smith et al."
]{biblatex}
\addbibresource{references.bib}  % Path relative to proposal.tex


% -----------------------------------------------------------------------------
% 10. HYPERLINKS & PDF METADATA
% hyperref must be loaded SECOND TO LAST (after almost all other packages).
% cleveref must be loaded LAST (after hyperref).
% -----------------------------------------------------------------------------
\usepackage[
  colorlinks = true,
  linkcolor  = blue,      % Internal links (sections, figures, tables)
  citecolor  = blue,      % Citation links
  urlcolor   = blue,      % External URLs
  pdftitle   = {PhD Proposal},
  pdfauthor  = {Your Name},
]{hyperref}

\usepackage{cleveref}     % Smart cross-references: \cref{fig:results} → "Figure 1"
                          % Also handles ranges: \crefrange{eq:1}{eq:3} → "Equations 1 to 3"
